\documentclass[../main.tex]{subfiles}
\begin{document}

\textit{Hệ thống được xây dựng với sự hỗ trợ của các agent lập trình AI. Chương này mô tả cách tổ
chức các agent để chúng làm việc như một nhà máy phần mềm nhỏ thay vì một trợ lý tự do: ngữ cảnh
được cung cấp, các trạm của dây chuyền, các rào chắn buộc agent tuân thủ quy tắc của dự án, các
mức kiểm thử phải vượt qua, và cách sản phẩm được đóng gói, vận hành và theo dõi.}

\section{Từ trợ lý đến nhà máy}
Một agent lập trình AI tự làm việc có hai điểm yếu. Nó không biết dự án nên phải đoán; và nó không
biết giới hạn của mình nên có thể ``sửa'' một quy tắc bảo mật dễ dàng như sửa lỗi chính tả. Một số
lỗi trong Bảng~\ref{tab:defects} cho thấy cả hai: lệnh triển khai suýt đưa cấu hình phát triển lên
production (F-37), và thay đổi đề xuất giả lập vai trò PM ngay trên trình duyệt (F-53). Lời giải gồm
bốn phần, mỗi phần là một mục dưới đây:
\begin{enumerate}
  \item \textbf{Kỹ thuật ngữ cảnh (context engineering)}---agent đọc đúng tài liệu trước khi hành động.
  \item \textbf{Nhà máy (factory)}---công việc đi qua các trạm riêng biệt với quyền hạn riêng.
  \item \textbf{Rào chắn (guardrails)}---quy tắc được thực thi bằng mã, không chỉ viết thành lời khuyên.
  \item \textbf{Kiểm chứng}---kiểm thử đơn vị, tích hợp, đầu-cuối, và kiểm thử chính các bài test.
\end{enumerate}

\begin{figure}[H]
\centering
\resizebox{\textwidth}{!}{\begin{tikzpicture}[
  st/.style={draw, rounded corners, fill=blue!6, minimum width=2.15cm, minimum height=1.25cm, align=center, font=\footnotesize},
  gate/.style={draw, fill=orange!15, minimum width=13.6cm, minimum height=0.7cm, font=\footnotesize, align=center},
  ctx/.style={draw, fill=green!8, minimum width=13.6cm, minimum height=0.7cm, font=\footnotesize, align=center},
  >=Stealth]
\node[ctx] (ctx) at (6.8,2.0) {Ngữ cảnh: CLAUDE.md + D1--D8 (yêu cầu, module, kiến trúc, hợp đồng, kiểm thử, chính sách, quyết định, RTM)};
\node[st] (s1) at (0.0,0) {1. Phân tích\\yêu cầu};
\node[st] (s2) at (2.27,0) {2. Kiến trúc\\sư};
\node[st] (s3) at (4.54,0) {3. Kỹ sư\\kiểm thử\\(đỏ)};
\node[st] (s4) at (6.81,0) {4. Lập trình\\(xanh)};
\node[st] (s5) at (9.08,0) {5. Kỹ sư KT\\(đột biến)};
\node[st] (s6) at (11.35,0) {6. Người\\review};
\node[st] (s7) at (13.62,0) {7. Kỹ sư\\phát hành};
\foreach \a/\b in {s1/s2,s2/s3,s3/s4,s4/s5,s5/s6,s6/s7} \draw[->] (\a) -- (\b);
\draw[->, dashed] (s6.south) to[bend left=25] node[below, font=\scriptsize]{có lỗi: quay lại lập trình} (s4.south);
\node[gate] (g) at (6.8,-2.3) {Rào chắn trên mọi lần sửa file và lệnh (cấm / hỏi người / cho phép) \quad + \quad cổng Definition of Done khi kết thúc};
\draw[->] (ctx.south) -- (6.8,0.75);
\draw[->] (g.north) -- (6.8,-1.6);
\end{tikzpicture}}
\caption{Nhà máy agent AI: ngữ cảnh đầu vào, các trạm nối tiếp, rào chắn quanh mọi bước}\label{fig:factory}
\end{figure}

\section{Kỹ thuật ngữ cảnh}
Kho mã có một lớp tài liệu điều hướng viết cho agent trước, cho người sau. Khi bắt đầu bất kỳ việc
gì, agent tự động nạp \code{CLAUDE.md}; tệp này liệt kê các lệnh kiểm chứng, các bất biến không được
phá vỡ, và thứ tự đọc tám tài liệu:

\begin{table}[H]
\centering
\caption{Các tài liệu agent đọc và câu hỏi mỗi tài liệu trả lời}\label{tab:context}
\renewcommand{\arraystretch}{1.2}
\begin{tabular}{|p{1.2cm}|p{4.2cm}|p{8.6cm}|}
\hline
\textbf{TL} & \textbf{Nội dung} & \textbf{Câu hỏi được trả lời cho agent} \\ \hline
D1 & Yêu cầu & Hệ thống phải làm gì? (mã FR, NFR) \\ \hline
D2 & Bản đồ module & File nào cài đặt việc này, cái gì phụ thuộc vào nó? \\ \hline
D3 & Kiến trúc & Thay đổi của tôi phải giữ quy tắc thứ tự hay phân lớp nào? \\ \hline
D4 & Hợp đồng API và dữ liệu & Đây có phải thay đổi hợp đồng? (nếu có: dừng và hỏi) \\ \hline
D5 & Kiểm thử và nghiệm thu & Điều gì chứng minh nó chạy đúng; thế nào là ``xong''? \\ \hline
D6 & Rủi ro, phân quyền, chính sách & Tôi được tự làm, làm rồi báo cáo, hay phải hỏi? \\ \hline
D7 & Sổ phát triển & Việc này đã được quyết định chưa, vì sao? Lỗi này từng gặp chưa? \\ \hline
D8 & Ma trận truy vết & File tôi đang sửa phục vụ yêu cầu và test nào? \\ \hline
\end{tabular}
\end{table}

Hai nguyên tắc giữ cho ngữ cảnh đáng tin. Thứ nhất, mọi tài liệu được viết \emph{dựa trên mã thật}
và cập nhật trong cùng commit với mã nó mô tả; bộ tài liệu cũ từng mô tả những hành vi mà mã không có,
và một agent tin vào tài liệu lỗi thời còn tệ hơn agent không có tài liệu. Thứ hai, mỗi trạm chỉ nạp
những gì cần (người phân tích đọc D1, D6, D8; kiến trúc sư đọc D2, D3, D4, D6, D7), giữ ngữ cảnh gọn
và tập trung.

\section{Nhà máy: các trạm và trách nhiệm}
Mỗi trạm là một định nghĩa agent riêng trong \code{.claude/agents/}, có chỉ dẫn, quyền dùng công cụ
và mô hình riêng. Lệnh \code{/feature <mô tả>} đưa một yêu cầu đi qua mọi trạm theo thứ tự;
\code{/verify} chạy mọi cổng chất lượng; \code{/release-check} đóng gói và chạy thử sản phẩm.

\begin{longtable}{|p{2.6cm}|p{6.4cm}|p{4.9cm}|}
\caption{Các trạm của nhà máy}\label{tab:stations}\\
\hline
\textbf{Trạm} & \textbf{Sản phẩm đầu ra} & \textbf{Không được phép} \\ \hline
\endfirsthead
\hline
\textbf{Trạm} & \textbf{Sản phẩm đầu ra} & \textbf{Không được phép} \\ \hline
\endhead
Phân tích yêu cầu & Yêu cầu có mã trong D1, tiêu chí nghiệm thu Given/When/Then trong D5, dòng mới trong D8. & Sửa mã nguồn. \\ \hline
Kiến trúc sư & Kế hoạch tác động: file kèm vùng tự chủ, ảnh hưởng hợp đồng, bất biến có rủi ro, danh sách test; một bản ghi quyết định trong D7. & Sửa mã nguồn; tiếp tục khi gặp mục vùng đỏ. \\ \hline
Kỹ sư kiểm thử & Viết test đơn vị và tích hợp \emph{trước} (phải thất bại); sau khi cài đặt, gieo một lỗi cho mỗi test để chứng minh test bắt được. & Nới lỏng hoặc xoá assertion cũ. \\ \hline
Lập trình & Thay đổi nhỏ nhất để test qua, rồi build production. & Sửa file ngoài kế hoạch. \\ \hline
Người review & PASS, hoặc danh sách lỗi kèm kịch bản: phân quyền, hợp đồng, bí mật, injection, chính sách ngoại tuyến, dịch thuật, test, tài liệu. & Sửa bất cứ thứ gì (chỉ đọc). \\ \hline
Kỹ sư phát hành & Bảng cổng: bộ test đơn vị và tích hợp, test rào chắn, build, stack cấu hình production, smoke test, health và metrics. & Triển khai hoặc push khi chưa có người duyệt. \\ \hline
\end{longtable}

Sự tách biệt chính là cơ chế kiểm soát: người review đã tự viết mã thì không còn là người review, và
người phân tích không sửa được mã thì không thể ``giải quyết'' yêu cầu bằng cách lặng lẽ đổi hành vi.

\section{Rào chắn và bảo mật}
Chỉ dẫn có thể bị bỏ qua; mã thì không. Vì vậy các vùng tự chủ trong D6 còn được cài đặt thành
\emph{hook} mà môi trường chạy agent gọi trước mọi lần sửa file, mọi lệnh shell, và khi agent định
kết thúc (Hình~\ref{fig:guard}). Chính sách là tập hàm thuần (\code{.claude/hooks/policy.py}); script
hook chỉ chuyển đổi giữa giao thức JSON của môi trường và chính sách. Nếu chính script lỗi, nó từ
chối---rào chắn hỏng không được biến thành cửa mở.

\begin{figure}[H]
\centering
\begin{tikzpicture}[
  b/.style={draw, rounded corners, minimum height=0.9cm, align=center, font=\footnotesize},
  >=Stealth]
\node[b, fill=blue!6, minimum width=2.6cm] (a) at (0,0) {Agent muốn\\sửa / chạy};
\node[b, fill=gray!10, minimum width=2.8cm] (h) at (4.0,0) {Hook\\\code{guard.py}};
\node[b, fill=gray!10, minimum width=2.8cm] (p) at (8.0,0) {Chính sách\\\code{policy.py}};
\node[b, fill=red!12, minimum width=2.6cm] (d) at (12.2,1.4) {CẤM\\(không bao giờ)};
\node[b, fill=orange!15, minimum width=2.6cm] (q) at (12.2,0) {HỎI\\(người quyết)};
\node[b, fill=green!12, minimum width=2.6cm] (o) at (12.2,-1.4) {CHO PHÉP\\(+ nhắc nhở)};
\draw[->] (a) -- node[above, font=\scriptsize]{JSON} (h);
\draw[->] (h) -- (p);
\draw[->] (p.east) -- (d.west); \draw[->] (p.east) -- (q.west); \draw[->] (p.east) -- (o.west);
\end{tikzpicture}
\caption{Mọi lần sửa file và mọi lệnh đều qua chính sách trước khi thực hiện}\label{fig:guard}
\end{figure}

\begin{longtable}{|p{1.9cm}|p{7.4cm}|p{4.6cm}|}
\caption{Một số quy tắc rào chắn}\label{tab:guardrules}\\
\hline
\textbf{Quyết định} & \textbf{Quy tắc} & \textbf{Lý do} \\ \hline
\endfirsthead
\hline
\textbf{Quyết định} & \textbf{Quy tắc} & \textbf{Lý do} \\ \hline
\endhead
Cấm & Mọi nội dung sửa có chứa API key, khoá riêng tư hoặc khoá ký JWT từng bị lộ. & Bí mật không bao giờ vào kho mã. \\ \hline
Cấm & Đọc hoặc in \code{.env}; sửa \code{.env}, thư mục nộp bài đã đóng băng hoặc thư mục sinh tự động. & Bí mật; sản phẩm bất biến. \\ \hline
Cấm & Sửa một migration đã commit. & Migration đã áp dụng là bất biến; phải thêm migration mới. \\ \hline
Cấm & \code{git push --force}; xoá đệ quy \code{/} hoặc \code{\textasciitilde}; lệnh triển khai tar-qua-ssh cũ. & Không thể khôi phục; F-37. \\ \hline
Hỏi & Sửa xác thực, mô hình dữ liệu, schema API, danh sách định dạng upload, cấu hình, CI hoặc chính rào chắn. & Vùng đỏ trong D6: bảo mật và hợp đồng. \\ \hline
Hỏi & \code{git push}, script triển khai, hạ cấp CSDL, xoá volume Docker. & Ảnh hưởng người khác hoặc xoá dữ liệu. \\ \hline
Cho phép + nhắc & Store, router, bản dịch, prompt AI, trang giới thiệu. & Vùng vàng: báo cáo và cập nhật tài liệu. \\ \hline
Chặn kết thúc & \code{source/} thay đổi mà không có tài liệu nào thay đổi, hoặc mã backend/frontend đổi mà không có test, hoặc model đổi mà không có migration. & Definition of Done (D5 \S4). \\ \hline
\end{longtable}

\textbf{Lần chặn đầu tiên là chính tác giả.} Vài phút sau khi bật hook, rào chắn đã từ chối một lệnh
do chính agent đang viết chương này gửi: một script \emph{ghi tài liệu} có nhắc tới lệnh triển khai
tar-qua-ssh cũ bị nhầm là đang chạy lệnh đó. Quy tắc được sửa để nội dung đang ghi vào file
(``heredoc'') được coi là dữ liệu, có thêm test cho đúng trường hợp này, và một test nữa kiểm tra lệnh
nguy hiểm đặt \emph{sau} đoạn nội dung đó vẫn bị chặn.

\textbf{Bảo mật trong chính sản phẩm.} Rào chắn bảo vệ kho mã; ứng dụng có lớp phòng thủ riêng, mô tả
ở Chương~3 và kiểm thử ở Chương~4: phân quyền phía server cho mọi route dự án, trả 404 thay vì 403
cho người ngoài, kiểm tra khi khởi động từ chối khoá yếu, CORS wildcard hay tài khoản demo trên
production, upload giới hạn theo danh sách cho phép với tên file do server sinh, và bình luận được
hiển thị dạng văn bản, không bao giờ dạng HTML.

\section{Kiểm thử đơn vị và kiểm thử tích hợp}
\begin{table}[H]
\centering
\caption{Các mức kiểm thử}\label{tab:levels}
\renewcommand{\arraystretch}{1.2}
\begin{tabular}{|p{2.6cm}|p{7.5cm}|p{2.0cm}|p{1.5cm}|}
\hline
\textbf{Mức} & \textbf{Nội dung kiểm thử} & \textbf{Công cụ} & \textbf{Số lượng} \\ \hline
Đơn vị (frontend) & Module thuần---sắp xếp phụ thuộc, đường găng, thứ tự khẩn cấp, bộ điều phối bàn phím, bộ phân tích mention---và hai store riêng lẻ. & Vitest & 442 \\ \hline
Đơn vị (rào chắn) & Mọi quy tắc chính sách, giao thức JSON của hook, hành vi từ chối khi lỗi. & pytest & 57 \\ \hline
Tích hợp (backend) & Router, phiên CSDL và xác thực thật qua HTTP test client; từng tầng AI giả lập ở biên HTTP; migration từ CSDL rỗng. & pytest & 185 \\ \hline
Tích hợp (giao diện) & Component gắn vào DOM giả lập cùng store. & Vitest & (trong 442) \\ \hline
Đầu-cuối & Instance công khai đã seed, mọi chức năng qua API, AI với mô hình thật. & script & 32 bước \\ \hline
\end{tabular}
\end{table}
Chính các bài test cũng được kiểm thử. Với rào chắn, bỏ quy tắc chặn force-push và quy tắc
migration làm năm test thất bại; khôi phục lại thì tất cả đạt. Tích hợp liên tục chạy ba job mỗi lần
push---frontend, backend và rào chắn---nên thay đổi làm hỏng bất kỳ phần nào cũng không thể lọt qua.

\section{Đóng gói, vận hành và theo dõi sản phẩm}
\subsection{Đóng gói}
Sản phẩm gồm bốn container mô tả bằng Docker Compose: MySQL, backend FastAPI, nginx phục vụ frontend
đã build, và \code{cloudflared} cho tunnel công khai. Image không chứa bí mật nào: \code{.dockerignore}
loại \code{.env} và CSDL, và script triển khai kiểm tra lại image đã build. Bí mật được sinh riêng cho
từng máy bằng \code{scripts/dev-env.sh}.

\subsection{Vận hành}
\code{scripts/local-prod.sh up} chạy migration ở một bước riêng, khởi động stack với cơ chế kiểm tra
an toàn production được bật, và xác minh tư thế (khoá không mặc định, không tài khoản demo, CORS cố
định, lược đồ ở phiên bản mới nhất). Docker tự khởi động lại container bị dừng, và file log được giới
hạn 5$\times$10\,MB mỗi container.

\subsection{Theo dõi}
Giai đoạn này bổ sung ba tín hiệu:
\begin{itemize}
  \item \textbf{Health có kiểm tra CSDL.} Trước đây \code{/api/health} vẫn báo ``healthy'' kể cả khi
        CSDL ngừng. Giờ nó chạy một truy vấn và trả HTTP 503 khi không kết nối được CSDL; healthcheck
        của Docker dùng endpoint này nên \code{docker ps} hiển thị đúng trạng thái.
  \item \textbf{Mã yêu cầu và log truy cập.} Mỗi phản hồi có \code{X-Request-ID}; mỗi yêu cầu ghi một
        dòng log gồm phương thức, route, mã trạng thái và thời gian. Log ghi \emph{mẫu} route
        (\code{/api/tasks/\{task\_id\}}) nên không id, token hay query string nào lọt vào log.
  \item \textbf{Metrics.} \code{/api/metrics} công bố theo định dạng Prometheus: số yêu cầu theo route
        và trạng thái, histogram độ trễ và \textbf{số câu trả lời AI theo tầng dự phòng}. Chỉ số cuối
        quan trọng nhất: khi nhà cung cấp AI lỗi, mọi endpoint vẫn trả lời (theo thiết kế), nên chỉ bộ
        đếm này cho biết người dùng đang nhận câu trả lời dự phòng thay vì mô hình. Trên production,
        metrics bị ẩn (404) trừ khi gửi kèm token bí mật.
\end{itemize}
\code{scripts/monitor.sh} gộp các tín hiệu thành một bản tóm tắt và cảnh báo khi có lỗi server hoặc
khi hơn một nửa câu trả lời AI đến từ tầng dự phòng. Kết quả trên instance công khai sau khi chạy
kiểm thử đầu-cuối:
\begin{verbatim}
health      HTTP 200  (0.0026s)
container   koshi-backend-prod-local   Up (healthy)
container   koshi-db-prod-local        Up 2 days (healthy)
requests    48 total, 0 server errors
AI answers  5, 0 from the offline fallback (0%)
\end{verbatim}
Tám test tích hợp bao phủ phần này, gồm: metrics bị ẩn khi thiếu token, mã yêu cầu độc hại bị thay
thế, và id của task không bao giờ xuất hiện trong nhãn metric.

\section{Đánh giá}
Nhà máy không làm agent trở nên không bao giờ sai. Cổng kết thúc chỉ kiểm tra tài liệu và test
\emph{có thay đổi}, không kiểm tra chúng đúng; việc đánh giá đó thuộc về trạm review và cuối cùng là
con người. Hook chỉ ràng buộc agent lập trình, vì vậy tích hợp liên tục chạy lại mọi bộ test một cách
độc lập. Trong giới hạn đó, những quy tắc trước đây phụ thuộc vào việc agent có chịu đọc tài liệu
chính sách hay không giờ được thực thi trên mọi hành động, và mỗi quy tắc đều có test đã được chứng
minh là thất bại khi quy tắc bị gỡ bỏ.

\end{document}
